\documentclass[11pt,a4paper]{article}
\usepackage{fontspec,amsmath,amssymb,amsthm,microtype}
\setmainfont{lmroman10-regular.otf}[BoldFont=lmroman10-bold.otf,ItalicFont=lmroman10-italic.otf,BoldItalicFont=lmroman10-bolditalic.otf]
\usepackage{xcolor}
\usepackage[most]{tcolorbox}
\usepackage{mathrsfs}
\providecommand{\lt}{<}
\providecommand{\gt}{>}
\setmonofont{DejaVuSansMono.ttf}[BoldFont=DejaVuSansMono-Bold.ttf,ItalicFont=DejaVuSansMono-Oblique.ttf,BoldItalicFont=DejaVuSansMono-BoldOblique.ttf,Scale=0.86]
\newfontfamily\leanfallback{STIXTwoMath-Regular.otf}[Scale=0.86]
\definecolor{leanrule}{gray}{0.62}
\definecolor{leanback}{gray}{0.962}
\definecolor{leantitle}{gray}{0.22}
\definecolor{leanmuted}{gray}{0.38}
\newtcolorbox{leanbox}[1]{enhanced,breakable,colback=leanback,colframe=leanrule,
  boxrule=0.4pt,arc=0.6mm,left=2.4mm,right=2.4mm,top=1.8mm,bottom=1.8mm,
  before skip=0.8em,after skip=1em,fonttitle=\small\bfseries,coltitle=leantitle,
  colbacktitle=leanback,titlerule=0.3pt,title={#1}}
\newlength{\leancharwidth}
\newenvironment{leancode}{\par\smallskip\begingroup\ttfamily\footnotesize\raggedright
  \settowidth{\leancharwidth}{M}\setlength{\parskip}{0pt}\setlength{\parindent}{0pt}
  \frenchspacing}{\par\endgroup\smallskip}
\newcommand{\leanline}[2]{\hangindent=\dimexpr#1\leancharwidth+4\leancharwidth\relax
  \hangafter=1\noindent\hspace*{#1\leancharwidth}#2\par}
\newcommand{\leansource}[1]{\par{\footnotesize\color{leanmuted}#1}\par}
\newcommand{\leanlabel}[1]{\par\smallskip\noindent{\small\bfseries #1}\par}
\usepackage[margin=25mm]{geometry}
\usepackage[unicode,hidelinks]{hyperref}
\hypersetup{pdfauthor={LangevinLean Project}}
\setlength{\emergencystretch}{3em}
\setlength{\parskip}{0.4em}
\setlength{\parindent}{0pt}
\allowdisplaybreaks[2]
\raggedbottom
\title{\texorpdfstring{Identify Gaussian proposals with their radial density}{Identify Gaussian proposals with their radial density}}
\author{LangevinLean Project}
\date{}
\begin{document}\maketitle
Draft mathematical exposition; these notes may contain mistakes.


For positive precision \(a\), let \(J(a)=\int e^{-a\|z\|^2}dz\) and \(C(a)=J(a)^{-1}\). Gaussian integrability and positivity give \(0<J(a)<\infty\). Translation then proves that \(q_{a,c}(z)=C(a)e^{-a\|z-c\|^2}\) integrates to one for every center \(c\).



The Gaussian Fourier integral gives characteristic function \(e^{-\|t\|^2/(4a)}e^{i\langle c,t\rangle}\). For \(a=1/(4h)\), it agrees with the law of \(c+\sqrt{2h}Z\). Equality of characteristic functions identifies the probability measures. Taking \(c=x-h\nabla U(x)\) gives the ULA proposal density, with covariance \(2hI\).


\section{\texorpdfstring{The radial Gaussian integral}{The radial Gaussian integral}}


Set \(J(a)=\int_H e^{-a\|z\|^2}dz\) as an ordinary real integral. The positive-precision results below prove its integrability and positivity before using \(J(a)^{-1}\) as a normalizing constant.



The radial integral is evaluated in the Euclidean volume of the state space. A positive precision gives a product of finite scalar Gaussian integrals, retaining its dimension dependence.

\begin{leanbox}{Lean: \texttt{radialGaussianIntegral}}
\leansource{\texttt{radialGaussianIntegral} \quad UnadjustedLangevin/\allowbreak{}GaussianProposalDensity.lean, lines 26–26}
\begin{leancode}
\leanline{0}{def\ radialGaussianIntegral\ (a\ :\ ℝ)\ :\ ℝ\ :=\ ∫\ z\ :\ H,\ Real.exp\ (-a\ *\ ‖z‖\ \textasciicircum{}\ 2)}
\end{leancode}
\leanlabel{How the Lean proof follows the argument}
\texttt{radial\allowbreak{}Gaussian\allowbreak{}Integral} is a definition with defining expression \texttt{\ensuremath{\int} z : H, Real.\allowbreak{}exp (-\allowbreak{}a * \ensuremath{\Vert}z\ensuremath{\Vert} \textasciicircum{} 2)}. The right-hand side specifies the object used in the mathematical statement.
\end{leanbox}

\section{\texorpdfstring{The radial normalization constant}{The radial normalization constant}}


Define \(C(a)=J(a)^{-1}\). The constant is defined by the volume integral on the given finite-dimensional Hilbert space and is independent of the center.



The normalizer is the reciprocal of the positive finite radial integral. This choice is fixed before translations of the kernel are considered.

\begin{leanbox}{Lean: \texttt{radialGaussianNormalizer}}
\leansource{\texttt{radialGaussianNormalizer} \quad UnadjustedLangevin/\allowbreak{}GaussianProposalDensity.lean, lines 29–29}
\begin{leancode}
\leanline{0}{def\ radialGaussianNormalizer\ (a\ :\ ℝ)\ :\ ℝ\ :=\ (radialGaussianIntegral\ (H\ :=\ H)\ a)⁻¹}
\end{leancode}
\leanlabel{How the Lean proof follows the argument}
\texttt{radial\allowbreak{}Gaussian\allowbreak{}Normalizer} is a definition with defining expression \texttt{H) a)\ensuremath{{}^{-1}}}. The right-hand side specifies the object used in the mathematical statement.
\end{leanbox}

\section{\texorpdfstring{Positivity of the radial integral}{Positivity of the radial integral}}


For \(a>0\), the radial Gaussian is integrable by the positive-precision Gaussian bound. The integrand is strictly positive, so its ordinary integral is strictly positive by positivity of the integral of an everywhere positive continuous weight.



The radial integrand is strictly positive and continuous. Its integral over any nonempty ball is positive, establishing a nonzero denominator for normalization.

\begin{leanbox}{Lean: \texttt{radialGaussianIntegral\_\allowbreak{}pos}}
\leansource{\texttt{radialGaussianIntegral\_\allowbreak{}pos} \quad UnadjustedLangevin/\allowbreak{}GaussianProposalDensity.lean, lines 31–33}
\begin{leancode}
\leanline{0}{theorem\ radialGaussianIntegral\_pos\ (a\ :\ ℝ)\ (ha\ :\ 0\ <\ a)\ :}
\leanline{4}{0\ <\ radialGaussianIntegral\ (H\ :=\ H)\ a}
\end{leancode}
\leanlabel{How the Lean proof follows the argument}
\texttt{radial\allowbreak{}Gaussian\allowbreak{}Integral\_\allowbreak{}pos}: the proof is the application \texttt{integral\_\allowbreak{}exp\_\allowbreak{}pos (integrable\_\allowbreak{}exp\_\allowbreak{}neg\_\allowbreak{}mul\_\allowbreak{}norm\_\allowbreak{}sq a ha)}. It specializes that result with the hypotheses and objects in the exact statement.
\end{leanbox}

\section{\texorpdfstring{Positivity of the normalizer}{Positivity of the normalizer}}


The inverse of a strictly positive real number is strictly positive. Apply this to \(J(a)\), whose positivity was proved immediately above. This also proves the normalizer is nonzero for later cancellation.



The reciprocal of a positive finite scalar is positive. Thus normalization preserves strict positivity of every Gaussian density value.

\begin{leanbox}{Lean: \texttt{radialGaussianNormalizer\_\allowbreak{}pos}}
\leansource{\texttt{radialGaussianNormalizer\_\allowbreak{}pos} \quad UnadjustedLangevin/\allowbreak{}GaussianProposalDensity.lean, lines 35–36}
\begin{leancode}
\leanline{0}{theorem\ radialGaussianNormalizer\_pos\ (a\ :\ ℝ)\ (ha\ :\ 0\ <\ a)\ :}
\leanline{4}{0\ <\ radialGaussianNormalizer\ (H\ :=\ H)\ a}
\end{leancode}
\leanlabel{How the Lean proof follows the argument}
\texttt{radial\allowbreak{}Gaussian\allowbreak{}Normalizer\_\allowbreak{}pos}: the proof is the application \texttt{inv\_\allowbreak{}pos.\allowbreak{}mpr (radial\allowbreak{}Gaussian\allowbreak{}Integral\_\allowbreak{}pos a ha)}. It specializes that result with the hypotheses and objects in the exact statement.
\end{leanbox}

\section{\texorpdfstring{The normalized radial density}{The normalized radial density}}


Define \(q_{a,c}(z)=C(a)e^{-a\|z-c\|^2}\), embedded in extended nonnegative scalars. The center changes the translated argument, while the normalizer depends only on precision and dimension.



The density is the radial exponential multiplied by its normalizing scalar. Translation to a nonzero mean later changes the exponent's center but leaves this scalar unchanged.

\begin{leanbox}{Lean: \texttt{radialGaussianDensity}}
\leansource{\texttt{radialGaussianDensity} \quad UnadjustedLangevin/\allowbreak{}GaussianProposalDensity.lean, lines 39–40}
\begin{leancode}
\leanline{0}{def\ radialGaussianDensity\ (a\ :\ ℝ)\ (c\ z\ :\ H)\ :\ ℝ≥0∞\ :=}
\leanline{2}{ENNReal.ofReal\ (radialGaussianNormalizer\ (H\ :=\ H)\ a\ *\ Real.exp\ (-a\ *\ ‖z\ -\ c‖\ \textasciicircum{}\ 2))}
\end{leancode}
\leanlabel{How the Lean proof follows the argument}
\texttt{radial\allowbreak{}Gaussian\allowbreak{}Density} is a definition with defining expression \texttt{H) a * Real.\allowbreak{}exp (-\allowbreak{}a * \ensuremath{\Vert}z -\allowbreak{} c\ensuremath{\Vert} \textasciicircum{} 2))}. The right-hand side specifies the object used in the mathematical statement.
\end{leanbox}

\section{\texorpdfstring{Integrability of the radial density}{Integrability of the radial density}}


Translation by \(c\) preserves Lebesgue measure, so \(e^{-a\|z-c\|^2}\) is integrable whenever \(a>0\). Multiply by the finite real constant \(C(a)\) to obtain integrability of the displayed real density.



The Gaussian exponential is integrable and the normalizer is finite. Scalar multiplication therefore preserves integrability of the normalized density.

\begin{leanbox}{Lean: \texttt{integrable\_\allowbreak{}radialGaussianDensity\_\allowbreak{}real}}
\leansource{\texttt{integrable\_\allowbreak{}radialGaussianDensity\_\allowbreak{}real} \quad UnadjustedLangevin/\allowbreak{}GaussianProposalDensity.lean, lines 42–45}
\begin{leancode}
\leanline{0}{theorem\ integrable\_radialGaussianDensity\_real\ (a\ :\ ℝ)\ (ha\ :\ 0\ <\ a)\ (c\ :\ H)\ :}
\leanline{4}{Integrable\ (fun\ z\ :\ H\ ↦\ radialGaussianNormalizer\ (H\ :=\ H)\ a\ *}
\leanline{6}{Real.exp\ (-a\ *\ ‖z\ -\ c‖\ \textasciicircum{}\ 2))\ volume}
\end{leancode}
\leanlabel{How the Lean proof follows the argument}
\texttt{integrable\_\allowbreak{}radial\allowbreak{}Gaussian\allowbreak{}Density\_\allowbreak{}real}: the source steps apply or rewrite \texttt{integrable\_\allowbreak{}exp\_\allowbreak{}neg\_\allowbreak{}mul\_\allowbreak{}norm\_\allowbreak{}sq}, \texttt{comp\_\allowbreak{}sub\_\allowbreak{}right}, \texttt{const\_\allowbreak{}mul}.
\end{leanbox}

\section{\texorpdfstring{Its unit integral}{Its unit integral}}


Translate \(z\mapsto z-c\) in the integral and pull out \(C(a)\). The result is \(C(a)J(a)=1\), justified by \(J(a)>0\). Thus normalization is independent of the center.



The chosen reciprocal cancels the radial integral, giving total integral one. Positivity and finiteness justify that cancellation in ordinary real arithmetic.

\begin{leanbox}{Lean: \texttt{integral\_\allowbreak{}radialGaussianDensity\_\allowbreak{}real}}
\leansource{\texttt{integral\_\allowbreak{}radialGaussianDensity\_\allowbreak{}real} \quad UnadjustedLangevin/\allowbreak{}GaussianProposalDensity.lean, lines 47–51}
\begin{leancode}
\leanline{0}{theorem\ integral\_radialGaussianDensity\_real\ (a\ :\ ℝ)\ (ha\ :\ 0\ <\ a)\ (c\ :\ H)\ :}
\leanline{4}{(∫\ z\ :\ H,\ radialGaussianNormalizer\ (H\ :=\ H)\ a\ *\ Real.exp\ (-a\ *\ ‖z\ -\ c‖\ \textasciicircum{}\ 2))\ =\ 1}
\end{leancode}
\leanlabel{How the Lean proof follows the argument}
\texttt{integral\_\allowbreak{}radial\allowbreak{}Gaussian\allowbreak{}Density\_\allowbreak{}real}: the source steps apply or rewrite \texttt{integral\_\allowbreak{}const\_\allowbreak{}mul}, \texttt{inv\_\allowbreak{}mul\_\allowbreak{}cancel\ensuremath{{}_{0}}}, \texttt{radial\allowbreak{}Gaussian\allowbreak{}Integral\_\allowbreak{}pos}.
\end{leanbox}

\section{\texorpdfstring{The radial Gaussian probability law}{The radial Gaussian probability law}}


The total mass of the weighted-density measure is the extended integral of its density. Nonnegativity and the proved real integrability identify it with a nonnegative scalar of the real integral, which is one. This constructs probability normalization, including finiteness.



A measurable nonnegative density of integral one defines a probability measure. The radial Gaussian density therefore specifies a probability measure.

\begin{leanbox}{Lean: \texttt{radialGaussianDensity\_\allowbreak{}isProbability}}
\leansource{\texttt{radialGaussianDensity\_\allowbreak{}isProbability} \quad UnadjustedLangevin/\allowbreak{}GaussianProposalDensity.lean, lines 54–63}
\begin{leancode}
\leanline{0}{theorem\ radialGaussianDensity\_isProbability\ (a\ :\ ℝ)\ (ha\ :\ 0\ <\ a)\ (c\ :\ H)\ :}
\leanline{4}{IsProbabilityMeasure\ (volume.withDensity\ (radialGaussianDensity\ a\ c))}
\end{leancode}
\leanlabel{How the Lean proof follows the argument}
\texttt{radial\allowbreak{}Gaussian\allowbreak{}Density\_\allowbreak{}is\allowbreak{}Probability}: the source steps apply or rewrite \texttt{with\allowbreak{}Density\_\allowbreak{}apply}, \texttt{Measurable\allowbreak{}Set.\allowbreak{}univ}, \texttt{set\allowbreak{}LIntegral\_\allowbreak{}univ}, \texttt{of\allowbreak{}Real\_\allowbreak{}integral\_\allowbreak{}eq\_\allowbreak{}lintegral\_\allowbreak{}of\allowbreak{}Real}.
\end{leanbox}

\section{\texorpdfstring{The complex-valued radial Gaussian integral}{The complex-valued radial Gaussian integral}}


Let \(a>0\) and write \(J(a)=\int e^{-a\|z\|^2}dz\), which is finite by the radial Gaussian integrability estimate. Identifying real numbers with their image in \(\mathbb C\), the pointwise identity \(\exp(-a\|z\|^2)=e^{-a\|z\|^2}\) identifies the complex exponential with the real Gaussian weight.



The real-to-complex inclusion is continuous and linear, so it commutes with the Bochner integral of this integrable function. Consequently the complex-valued radial integral equals the complex image of \(J(a)\). This is the normalization used in the subsequent Fourier calculation.

\begin{leanbox}{Lean: \texttt{integral\_\allowbreak{}complex\_\allowbreak{}radialGaussian}}
\leansource{\texttt{integral\_\allowbreak{}complex\_\allowbreak{}radialGaussian} \quad UnadjustedLangevin/\allowbreak{}GaussianProposalDensity.lean, lines 66–78}
\begin{leancode}
\leanline{0}{theorem\ integral\_complex\_radialGaussian\ (a\ :\ ℝ)\ (\_ha\ :\ 0\ <\ a)\ :}
\leanline{4}{(∫\ z\ :\ H,\ Complex.exp\ (-(a\ :\ ℂ)\ *\ ‖z‖\ \textasciicircum{}\ 2))\ =}
\leanline{6}{(radialGaussianIntegral\ (H\ :=\ H)\ a\ :\ ℂ)}
\end{leancode}
\leanlabel{How the Lean proof follows the argument}
\texttt{integral\_\allowbreak{}complex\_\allowbreak{}radial\allowbreak{}Gaussian}: the source steps apply or rewrite \texttt{integral\_\allowbreak{}congr\_\allowbreak{}ae}, \texttt{ae\_\allowbreak{}of\_\allowbreak{}all}, \texttt{Complex.\allowbreak{}of\allowbreak{}Real\_\allowbreak{}exp}. The \texttt{calc} block checks the successive equalities or inequalities in the argument. The remaining scalar rearrangements are checked by \texttt{ring}.
\end{leanbox}

\section{\texorpdfstring{The Gaussian character integral}{The Gaussian character integral}}


Apply the finite-dimensional complex Gaussian Fourier integral with imaginary linear coefficient. Combining the two exponentials and using \(i^2=-1\) gives \(\int e^{-a\|z\|^2}e^{i\langle z,t\rangle}dz=J(a)e^{-\|t\|^2/(4a)}\). The zero-frequency Gaussian formula identifies the normalization factor with the same real integral \(J(a)\).



The character is \(e^{i\langle v,x\rangle}\). Inserting it into the normalized radial integral gives the characteristic expression with the corresponding isotropic covariance.

\begin{leanbox}{Lean: \texttt{integral\_\allowbreak{}radialGaussian\_\allowbreak{}character}}
\leansource{\texttt{integral\_\allowbreak{}radialGaussian\_\allowbreak{}character} \quad UnadjustedLangevin/\allowbreak{}GaussianProposalDensity.lean, lines 81–106}
\begin{leancode}
\leanline{0}{theorem\ integral\_radialGaussian\_character\ (a\ :\ ℝ)\ (ha\ :\ 0\ <\ a)\ (t\ :\ H)\ :}
\leanline{4}{(∫\ z\ :\ H,\ (Real.exp\ (-a\ *\ ‖z‖\ \textasciicircum{}\ 2)\ :\ ℂ)\ *}
\leanline{6}{Complex.exp\ ((inner\ ℝ\ z\ t\ :\ ℝ)\ *\ Complex.I))\ =}
\leanline{6}{(radialGaussianIntegral\ (H\ :=\ H)\ a\ :\ ℂ)\ *}
\leanline{8}{Complex.exp\ (-(‖t‖\ \textasciicircum{}\ 2\ :\ ℝ)\ /\ (4\ *\ (a\ :\ ℂ)))}
\end{leancode}
\leanlabel{How the Lean proof follows the argument}
\texttt{integral\_\allowbreak{}radial\allowbreak{}Gaussian\_\allowbreak{}character}: the intermediate facts include \texttt{h}, \texttt{hzero}, in source order; the source steps apply or rewrite \texttt{Complex.\allowbreak{}I}, \texttt{integral\_\allowbreak{}complex\_\allowbreak{}radial\allowbreak{}Gaussian}, \texttt{Complex.\allowbreak{}of\allowbreak{}Real\_\allowbreak{}neg}, \texttt{Complex.\allowbreak{}of\allowbreak{}Real\_\allowbreak{}pow}, \texttt{Complex.\allowbreak{}I\_\allowbreak{}sq}. The remaining scalar rearrangements are checked by \texttt{ring}.
\end{leanbox}

\section{\texorpdfstring{The characteristic function of the radial law}{The characteristic function of the radial law}}


Change from the density measure to a weighted Lebesgue integral. Translate \(z=u+c\); the character splits into \(e^{i\langle u,t\rangle}e^{i\langle c,t\rangle}\). The Fourier formula gives a factor \(J(a)\), canceled by \(C(a)\), leaving \(e^{-\|t\|^2/(4a)}e^{i\langle c,t\rangle}\).



The computed characteristic function matches that of the Gaussian with the displayed mean and covariance. Uniqueness of characteristic functions then identifies the radial density law with that Gaussian measure.

\begin{leanbox}{Lean: \texttt{charFun\_\allowbreak{}radialGaussianDensity}}
\leansource{\texttt{charFun\_\allowbreak{}radialGaussianDensity} \quad UnadjustedLangevin/\allowbreak{}GaussianProposalDensity.lean, lines 109–142}
\begin{leancode}
\leanline{0}{theorem\ charFun\_radialGaussianDensity\ (a\ :\ ℝ)\ (ha\ :\ 0\ <\ a)\ (c\ t\ :\ H)\ :}
\leanline{4}{charFun\ (volume.withDensity\ (radialGaussianDensity\ a\ c))\ t\ =}
\leanline{6}{Complex.exp\ (-(‖t‖\ \textasciicircum{}\ 2\ :\ ℝ)\ /\ (4\ *\ (a\ :\ ℂ)))\ *}
\leanline{8}{Complex.exp\ ((inner\ ℝ\ c\ t\ :\ ℝ)\ *\ Complex.I)}
\end{leancode}
\leanlabel{How the Lean proof follows the argument}
\texttt{char\allowbreak{}Fun\_\allowbreak{}radial\allowbreak{}Gaussian\allowbreak{}Density}: the intermediate facts include \texttt{hm}, \texttt{htrans}, in source order; the source steps apply or rewrite \texttt{char\allowbreak{}Fun\_\allowbreak{}apply}, \texttt{integral\_\allowbreak{}with\allowbreak{}Density\_\allowbreak{}eq\_\allowbreak{}integral\_\allowbreak{}to\allowbreak{}Real\_\allowbreak{}smul}, \texttt{integral\_\allowbreak{}mul\_\allowbreak{}const}, \texttt{integral\_\allowbreak{}radial\allowbreak{}Gaussian\_\allowbreak{}character}, \texttt{Complex.\allowbreak{}of\allowbreak{}Real\_\allowbreak{}inv}. The \texttt{calc} block checks the successive equalities or inequalities in the argument. The remaining scalar rearrangements are checked by \texttt{ring}.
\end{leanbox}

\section{\texorpdfstring{The proposal normalization constant}{The proposal normalization constant}}


Specialize the precision to \(a=1/(4h)\). The resulting constant normalizes a Gaussian with covariance \(2hI\). Its independence of the current state allows the same factor to cancel in the forward and reverse Metropolis ratio.



For proposal variance \(2h\), the normalizer becomes \((4\pi h)^{-d/2}\). Positive \(h\) is necessary for both this power and the variance interpretation.

\begin{leanbox}{Lean: \texttt{gaussianProposalNormalizer}}
\leansource{\texttt{gaussianProposalNormalizer} \quad UnadjustedLangevin/\allowbreak{}GaussianProposalDensity.lean, lines 145–146}
\begin{leancode}
\leanline{0}{def\ gaussianProposalNormalizer\ (h\ :\ ℝ)\ :\ ℝ\ :=}
\leanline{2}{radialGaussianNormalizer\ (H\ :=\ H)\ (1\ /\ (4\ *\ h))}
\end{leancode}
\leanlabel{How the Lean proof follows the argument}
\texttt{gaussian\allowbreak{}Proposal\allowbreak{}Normalizer} is a definition with defining expression \texttt{H) (1 /\allowbreak{} (4 * h))}. The right-hand side specifies the object used in the mathematical statement.
\end{leanbox}

\section{\texorpdfstring{The Gaussian proposal density}{The Gaussian proposal density}}


Define \(q_h(c,y)=C(1/(4h))e^{-\|y-c\|^2/(4h)}\). For \(h>0\), this is the density of a covariance-\(2hI\) Gaussian centered at \(c\). The density identification is proved later rather than built into the definition.



The proposal exponent is the squared distance from its deterministic center divided by \(4h\). Its density is therefore a translation of the normalized radial Gaussian.

\begin{leanbox}{Lean: \texttt{gaussianProposalDensity}}
\leansource{\texttt{gaussianProposalDensity} \quad UnadjustedLangevin/\allowbreak{}GaussianProposalDensity.lean, lines 149–151}
\begin{leancode}
\leanline{0}{def\ gaussianProposalDensity\ (h\ :\ ℝ)\ (c\ y\ :\ H)\ :\ ℝ≥0∞\ :=}
\leanline{2}{ENNReal.ofReal\ (gaussianProposalNormalizer\ (H\ :=\ H)\ h\ *}
\leanline{4}{Real.exp\ (-‖y\ -\ c‖\ \textasciicircum{}\ 2\ /\ (4\ *\ h)))}
\end{leancode}
\leanlabel{How the Lean proof follows the argument}
\texttt{gaussian\allowbreak{}Proposal\allowbreak{}Density} is a definition with defining expression \texttt{H) h *}. The right-hand side specifies the object used in the mathematical statement.
\end{leanbox}

\section{\texorpdfstring{Positivity of the proposal normalizer}{Positivity of the proposal normalizer}}


For \(h>0\), the precision \(1/(4h)\) is positive. Apply positivity of the radial Gaussian normalizer at that precision. This nonzero constant is used in detailed-balance cancellation.



Positive time makes the Gaussian normalizer strictly positive. Combined with the positive exponential, this yields positivity of the whole proposal kernel.

\begin{leanbox}{Lean: \texttt{gaussianProposalNormalizer\_\allowbreak{}pos}}
\leansource{\texttt{gaussianProposalNormalizer\_\allowbreak{}pos} \quad UnadjustedLangevin/\allowbreak{}GaussianProposalDensity.lean, lines 153–154}
\begin{leancode}
\leanline{0}{theorem\ gaussianProposalNormalizer\_pos\ (h\ :\ ℝ)\ (hh\ :\ 0\ <\ h)\ :}
\leanline{4}{0\ <\ gaussianProposalNormalizer\ (H\ :=\ H)\ h}
\end{leancode}
\leanlabel{How the Lean proof follows the argument}
\texttt{gaussian\allowbreak{}Proposal\allowbreak{}Normalizer\_\allowbreak{}pos}: the proof is the application \texttt{radial\allowbreak{}Gaussian\allowbreak{}Normalizer\_\allowbreak{}pos \_\allowbreak{} (by positivity)}. It specializes that result with the hypotheses and objects in the exact statement.
\end{leanbox}

\section{\texorpdfstring{Equality of radial and proposal density formulas}{Equality of radial and proposal density formulas}}


Unfold the two density definitions and the proposal normalizer. Their exponents agree by \(-(1/(4h))\|y-c\|^2=-\|y-c\|^2/(4h)\). Function extensionality gives equality of the density functions, with no probabilistic argument needed.



Substitution of the translated center identifies the two density formulas pointwise. The equality retains the same Euclidean norm and the same normalization constant.

\begin{leanbox}{Lean: \texttt{radialGaussianDensity\_\allowbreak{}eq\_\allowbreak{}proposalDensity}}
\leansource{\texttt{radialGaussianDensity\_\allowbreak{}eq\_\allowbreak{}proposalDensity} \quad UnadjustedLangevin/\allowbreak{}GaussianProposalDensity.lean, lines 156–162}
\begin{leancode}
\leanline{0}{theorem\ radialGaussianDensity\_eq\_proposalDensity\ (h\ :\ ℝ)\ (c\ :\ H)\ :}
\leanline{4}{radialGaussianDensity\ (1\ /\ (4\ *\ h))\ c\ =\ gaussianProposalDensity\ h\ c}
\end{leancode}
\leanlabel{How the Lean proof follows the argument}
\texttt{radial\allowbreak{}Gaussian\allowbreak{}Density\_\allowbreak{}eq\_\allowbreak{}proposal\allowbreak{}Density}: the proof is the application \texttt{funext y}. It specializes that result with the hypotheses and objects in the exact statement.
\end{leanbox}

\section{\texorpdfstring{The density of an affine Gaussian proposal}{The density of an affine Gaussian proposal}}


Compute characteristic functions on both sides. The affine image \(c+\sqrt{2h}Z\) has character \(e^{-h\|t\|^2}e^{i\langle c,t\rangle}\); the radial-density formula with \(a=1/(4h)\) gives the same expression. Both measures are probabilities, so uniqueness of a probability measure from its characteristic function identifies them. The square-root identity uses \(h>0\).



The affine noise map has Gaussian mean equal to the proposal center and covariance \(2hI_d\). Its characteristic function agrees with that of the explicit proposal density.

\begin{leanbox}{Lean: \texttt{gaussian\_\allowbreak{}proposal\_\allowbreak{}density}}
\leansource{\texttt{gaussian\_\allowbreak{}proposal\_\allowbreak{}density} \quad UnadjustedLangevin/\allowbreak{}GaussianProposalDensity.lean, lines 166–187}
\begin{leancode}
\leanline{0}{theorem\ gaussian\_proposal\_density\ (h\ :\ ℝ)\ (hh\ :\ 0\ <\ h)\ (c\ :\ H)\ :}
\leanline{4}{(stdGaussian\ H).map\ (fun\ ξ\ ↦\ c\ +\ Real.sqrt\ (2\ *\ h)\ •\ ξ)\ =}
\leanline{6}{volume.withDensity\ (gaussianProposalDensity\ h\ c)}
\end{leancode}
\leanlabel{How the Lean proof follows the argument}
\texttt{gaussian\_\allowbreak{}proposal\_\allowbreak{}density}: the intermediate facts include \texttt{ha}, \texttt{hm}, \texttt{hnorm}, \texttt{hnorm\allowbreak{}C}, in source order; the source steps apply or rewrite \texttt{radial\allowbreak{}Gaussian\allowbreak{}Density\_\allowbreak{}eq\_\allowbreak{}proposal\allowbreak{}Density}, \texttt{Measure.\allowbreak{}ext\_\allowbreak{}of\_\allowbreak{}char\allowbreak{}Fun}, \texttt{Real.\allowbreak{}norm\_\allowbreak{}eq\_\allowbreak{}abs}, \texttt{sq\_\allowbreak{}abs}, \texttt{Real.\allowbreak{}sq\_\allowbreak{}sqrt}. The remaining scalar rearrangements are checked by \texttt{field\_\allowbreak{}simp}.
\end{leanbox}

\section{\texorpdfstring{The radial representation of the proposal law}{The radial representation of the proposal law}}


Rewrite the proposal density using the equality between the translated radial Gaussian and proposal density formulas in the established affine-Gaussian law equality. This yields the equivalent precision-parameter form used in the detailed-balance proof.



Rewriting the proposal as a translated radial law separates the shape of the Gaussian from its center. The normalizer in this form is independent of the starting state.

\begin{leanbox}{Lean: \texttt{gaussian\_\allowbreak{}proposal\_\allowbreak{}radial\_\allowbreak{}density}}
\leansource{\texttt{gaussian\_\allowbreak{}proposal\_\allowbreak{}radial\_\allowbreak{}density} \quad UnadjustedLangevin/\allowbreak{}GaussianProposalDensity.lean, lines 190–194}
\begin{leancode}
\leanline{0}{theorem\ gaussian\_proposal\_radial\_density\ (h\ :\ ℝ)\ (hh\ :\ 0\ <\ h)\ (c\ :\ H)\ :}
\leanline{4}{(stdGaussian\ H).map\ (fun\ ξ\ ↦\ c\ +\ Real.sqrt\ (2\ *\ h)\ •\ ξ)\ =}
\leanline{6}{volume.withDensity\ (radialGaussianDensity\ (1\ /\ (4\ *\ h))\ c)}
\end{leancode}
\leanlabel{How the Lean proof follows the argument}
\texttt{gaussian\_\allowbreak{}proposal\_\allowbreak{}radial\_\allowbreak{}density}: the source steps apply or rewrite \texttt{radial\allowbreak{}Gaussian\allowbreak{}Density\_\allowbreak{}eq\_\allowbreak{}proposal\allowbreak{}Density}, \texttt{gaussian\_\allowbreak{}proposal\_\allowbreak{}density}.
\end{leanbox}

\section{\texorpdfstring{The radial representation of the ULA proposal}{The radial representation of the ULA proposal}}


Set the affine Gaussian center to \(c=x-h\nabla U(x)\). The fixed-state density theorem then gives \(x-h\nabla U(x)+\sqrt{2h}Z\) its exact radial density. No regularity of \(U\) is required because \(x\) is fixed throughout this statement.



For ULA the center is \(x-h\nabla U(x)\). Substituting this center into the radial representation gives the update law, with its drift dependence entirely in the translation.

\begin{leanbox}{Lean: \texttt{ula\_\allowbreak{}proposal\_\allowbreak{}radial\_\allowbreak{}density}}
\leansource{\texttt{ula\_\allowbreak{}proposal\_\allowbreak{}radial\_\allowbreak{}density} \quad UnadjustedLangevin/\allowbreak{}GaussianProposalDensity.lean, lines 197–200}
\begin{leancode}
\leanline{0}{theorem\ ula\_proposal\_radial\_density\ (U\ :\ H\ →\ ℝ)\ (h\ :\ ℝ)\ (hh\ :\ 0\ <\ h)\ (x\ :\ H)\ :}
\leanline{4}{(stdGaussian\ H).map\ (fun\ ξ\ ↦\ x\ -\ h\ •\ gradient\ U\ x\ +\ Real.sqrt\ (2\ *\ h)\ •\ ξ)\ =}
\leanline{6}{volume.withDensity\ (radialGaussianDensity\ (1\ /\ (4\ *\ h))\ (x\ -\ h\ •\ gradient\ U\ x))}
\end{leancode}
\leanlabel{How the Lean proof follows the argument}
\texttt{ula\_\allowbreak{}proposal\_\allowbreak{}radial\_\allowbreak{}density}: the proof is the application \texttt{gaussian\_\allowbreak{}proposal\_\allowbreak{}radial\_\allowbreak{}density h hh \_\allowbreak{}}. It specializes that result with the hypotheses and objects in the exact statement.
\end{leanbox}

\section{\texorpdfstring{The density of the ULA transition}{The density of the ULA transition}}


Expand the radial density at center \(x-h\nabla U(x)\). The exponent becomes \(-\|y-x+h\nabla U(x)\|^2/(4h)\), with the same positive normalizer. This is the explicit ULA transition density used in the Metropolis acceptance ratio.



The explicit density formula therefore describes the same transition as the independent Gaussian update. This law identification supports both mixture-density calculations and Metropolis proposal-ratio computations.

\begin{leanbox}{Lean: \texttt{ula\_\allowbreak{}proposal\_\allowbreak{}density}}
\leansource{\texttt{ula\_\allowbreak{}proposal\_\allowbreak{}density} \quad UnadjustedLangevin/\allowbreak{}GaussianProposalDensity.lean, lines 204–213}
\begin{leancode}
\leanline{0}{theorem\ ula\_proposal\_density\ (U\ :\ H\ →\ ℝ)\ (h\ :\ ℝ)\ (hh\ :\ 0\ <\ h)\ (x\ :\ H)\ :}
\leanline{4}{(stdGaussian\ H).map\ (fun\ ξ\ ↦\ x\ -\ h\ •\ gradient\ U\ x\ +\ Real.sqrt\ (2\ *\ h)\ •\ ξ)\ =}
\leanline{6}{volume.withDensity\ (fun\ y\ ↦\ ENNReal.ofReal\ (gaussianProposalNormalizer\ (H\ :=\ H)\ h\ *}
\leanline{8}{Real.exp\ (-‖y\ -\ x\ +\ h\ •\ gradient\ U\ x‖\ \textasciicircum{}\ 2\ /\ (4\ *\ h))))}
\end{leancode}
\leanlabel{How the Lean proof follows the argument}
\texttt{ula\_\allowbreak{}proposal\_\allowbreak{}density}: the intermediate fact is \texttt{heq}, in source order; the source steps apply or rewrite \texttt{gaussian\_\allowbreak{}proposal\_\allowbreak{}density}.
\end{leanbox}

\end{document}
